\chapter{Introducción}
\label{ch:intro}

Los polinomios ortogonales constituyen un área fundamental de la matemática con numerosas aplicaciones. En este trabajo nos centramos en una generalización particular: los polinomios ortogonales de Sobolev. La característica distintiva de estos polinomios es que su ortogonalidad viene dada por productos internos que involucran derivadas, lo que resulta en propiedades notablemente diferentes a las de los polinomios ortogonales clásicos.

\section{Definiciones}

Antes de adentrarnos en el estudio de los ceros de los polinomios ortogonales de Sobolev, es esencial establecer el marco teórico que usaremos a lo largo del trabajo. En este sentido, comenzaremos por definir los siguientes conceptos del álgebra lineal.

\begin{definition}\label{def:hermitian_matrix}
	\cite{Horn_Johnson_2012} Una matriz $A \in \mathcal{M}_{n}(\mathbb{C})$ se denomina \textit{hermitiana} si $A = A^*$, donde $A^*$ denota la matriz conjugada transpuesta de $A$. En el caso real, una matriz simétrica satisface $A = A^t$.

	Una matriz hermitiana $A$ se dice:
	\begin{enumerate}
		\item \textit{Definida positiva} (HPD) si para todo vector no nulo $x \in \mathbb{C}^n$, se cumple $x^*Ax > 0$. Equivalentemente, si todos sus autovalores son estrictamente positivos.

		\item \textit{Semidefinida positiva} (HPSD) si para todo vector $x \in \mathbb{C}^n$, se cumple $x^*Ax \geq 0$. Equivalentemente, si todos sus autovalores son no negativos.
	\end{enumerate}

	Para una matriz hermitiana, el producto escalar $\langle x, y \rangle_A = x^*Ay$ define un producto interno si y solo si $A$ es definida positiva.
\end{definition}

\begin{definition}\label{def:spectral_radius}
	\cite{dunford1963linear} El \textit{radio espectral} de una matriz cuadrada $A \in \mathcal{M}_{n}(\mathbb{C})$, denotado por $\rho(A)$, se define como el máximo de los valores absolutos de sus autovalores:
	\begin{equation}
		\rho(A) = \max\{|\lambda| : \lambda \text{ es autovalor de } A\}
	\end{equation}
	donde los autovalores son los escalares $\lambda \in \mathbb{C}$ tales que existe un vector no nulo $v \in \mathbb{C}^n$ que satisface $Av = \lambda v$.
\end{definition}

\begin{definition}\label{def:singular_values}
	Sea $A \in \mathcal{M}_{n}(\mathbb{C})$ matriz cuadrada,  la matriz $\mathbf{A}^{*}\mathbf{A}$  ( $\mathbf{A}^{t}\mathbf{A}$ en el caso de matrices con entradas reales) es una matriz hermitiana y semidefinida positiva. Por lo tanto, tiene $n$ autovalores mayores o iguales que cero contados con su multiplicidad
	$$
		0 \leq \lambda_{n,n}\leq \dots \leq \lambda_{2,n} \leq \lambda_{1,n}
	$$


	\noindent Los valores singulares de $\mathbf{A}$ son:
	$$
		\sigma_{k,n}= \sqrt{\lambda_{k,n}}
	$$
	\noindent siendo:
	$$
		0 \leq \sigma_{n,n} \leq \dots \sigma_{2,n} \leq \sigma_{1,n}
	$$

	\noindent Utilizando el criterio de Rayleigh y el Principio de min-max de  Courant-Fischer \cite{gallier2017spectral} se tiene que estos valores singulares tienen relación con la norma de aplicaciones lineales y continuas cuando consideramos la norma euclídea en los espacios $\mathbb{C}^n$. De hecho,
	\begin{enumerate}
		\item $$
			      \sigma_{1,n}^2 =
			      \max \{ <\mathbf{A}^{t}\mathbf{A}x,x>: \Vert x \Vert =1 \}
			      = \max \{\Vert \mathbf{A} x \Vert^2 : \Vert x \Vert =1\}=
			      \Vert \mathbf{A} \Vert^2
		      $$
		\item Principio min-max: el segundo mayor valor singular de una matriz es el mínimo de las normas de las restricciones de la aplicación $\mathbf{A}$ a los espacios de codimnesión $1$.
		      $$
			      \sigma_{2,n}^2 = \min_{V: cod(V)=1} \max \{\Vert \mathbf{A}x \Vert^2 , x\in V, \Vert x \Vert =1\}= \min_{V: cod(V)=1} \Vert \mathbf{A}/V \Vert^2
		      $$

		      \noindent siendo $\mathbf{A}/V$ la restricción de la aplicación lineal $\mathbf{A}$ al subespacio $V$.

	\end{enumerate}
\end{definition}

\noindent Una vez establecidos los conceptos básicos del álgebra lineal, debemos extender nuestro marco teórico al contexto de espacios funcionales, donde se desarrolla el estudio de los polinomios ortogonales de Sobolev. Para ello, necesitamos introducir las nociones de espacios prehilbertianos y espacios de Hilbert, que nos permitirán formalizar los distintos productos internos respecto a los cuales definiremos la ortogonalidad. Estos productos internos pueden tomar diversas formas, ya sean continuos o discretos, reales o complejos, y son fundamentales para caracterizar las propiedades de los polinomios ortogonales asociados.

\begin{definition}\label{def:prehilbert_space}
	Un \textit{espacio prehilbertiano} o \textit{espacio con producto interno} $(X, \langle \cdot, \cdot \rangle)$ es un espacio vectorial $X$ dotado de un producto interno $\langle \cdot, \cdot \rangle: X \times X \rightarrow \mathbb{C}$ (o $\mathbb{R}$), que satisface las siguientes propiedades para todo $x, y, z \in X$ y $\alpha \in \mathbb{C}$ (o $\mathbb{R}$):
	\begin{enumerate}
		\item $\langle x, y \rangle = \overline{\langle y, x \rangle}$ (simetría hermitiana)
		\item $\langle \alpha x + \beta y, z \rangle = \alpha \langle x, z \rangle + \beta \langle y, z \rangle$ (linealidad en la primera variable)
		\item $\langle x, x \rangle \geq 0$ y $\langle x, x \rangle = 0$ si y solo si $x = 0$ (definida positiva)
	\end{enumerate}

	Un \textit{espacio de Hilbert} es un espacio prehilbertiano que además es completo respecto a la norma inducida por el producto interno $\|x\| = \sqrt{\langle x, x \rangle}$, es decir, toda sucesión de Cauchy converge a un elemento del espacio.
\end{definition}

\begin{definition}\label{def:circum_lebesgue_measure}
	Sea
	\(
	S_{a}(r)\;=\;\bigl\{\,z\in\mathbb{C}\;:\;|z-a|=r\bigr\}
	\)
	la circunferencia de centro \(a\in\mathbb{C}\) y radio \(r>0\).
	Denotaremos por \(dm_{a;r}\) la \textbf{medida de Lebesgue} en \(S_{a}(r)\).
	Para toda función continua \(f:S_{a}(r)\to\mathbb{C}\) definimos
	$$
		\int_{S_{a}(r)} f(z)\,dm_{a;r}(z)
		\;:=\;
		\frac{1}{2\pi}\int_{0}^{2\pi} f\!\bigl(a+re^{i\theta}\bigr)\,d\theta,
	$$
	de modo que \(dm_{a;r}\) queda normalizada
	\(\bigl(dm_{a;r}(S_{a}(r))=1\bigr)\).
\end{definition}

Los productos internos que se trabajarán son de la forma:
\begin{enumerate}
	\item Continuo Real con Soportes Compactos
	      \begin{equation}
		      \langle p(t), q(t) \rangle = \int_{a}^{b} p(t) q(t) w_1(t) \, dt + \int_{c}^{d} p'(t) q'(t) w_2(t) \, dt
	      \end{equation}
	      donde $w_1(t)$ y $w_2(t)$ son pesos continuos, no negativos y distintos de la función $f \equiv 0$.
	\item Discreto Real con Átomo
	      \begin{equation}
		      \langle p(t), q(t) \rangle = \int_{a}^{b} p(t) q(t) w_1(t) \, dt + p'(c) q'(c)
	      \end{equation}
	      con $w_1(t)$ peso continuo, no negativo y distinto de la función $f \equiv 0$.
	\item Continuo Complejo en Circunferencias
	      \begin{equation}
		      \langle p(z) , q(z) \rangle_{S}\;=\;
		      \int_{S_{b}(r_b)}     p(z)\,\overline{q(z)} \, dm_{b;r_b}
		      \;+\;
		      \int_{S_{a}(r_a)}
		      p'(z)\,\overline{q'(z)} \, dm_{a;r_a}.
	      \end{equation}
	\item Discreto Complejo en Circunferencia con Átomo
	      \begin{equation}
		      \langle p(z) , q(z) \rangle_{S}\;=\;
		      \int_{S_{b}(r_b)} p(z)\,\overline{q(z)}\,dm_{b;r_b}
		      \;+\;
		      p'(c)\,\overline{q'(c)}.
	      \end{equation}
\end{enumerate}

Desde una perspectiva matricial, que facilita el análisis algebraico, podemos considerar que una matriz infinita HPD $M$ induce en $\mathbb{P}[z]$ el siguiente producto:
\begin{equation}
	\langle p(z), q(z) \rangle_M = v M w^{*},
\end{equation}
donde $v = (v_0, v_1, \ldots, v_n, 0, \ldots)$, $w = (w_0, w_1, \ldots, w_m, 0, \ldots) \in c_{00} = \{(x_n)_{n=1}^{\infty} \in \mathbb{C}^{\mathbb{N}} : \exists n_0 \in \mathbb{N} \quad | \quad x_n=0 ,\forall n\geq n_0\}$ para polinomios $p(z) = \sum_{k=0}^n v_k z^k$ y $q(z) = \sum_{k=0}^m w_k z^k$. La matriz de momentos $M = (c_{ij})^{\infty}_{i,j=0}$ definida por: $c_{ij} = \langle z^i, z^j \rangle_{M}$ Esto induce una norma Hilbertiana denotada por: $\|p(z)\|^{2}_{M} = \langle p(z), p(z)\rangle_{M},\forall p(z) \in P[z]$

\begin{definition}\label{def:sobolev_matrix_product}
	\cite{escribano2025boundedness} Sea $\{\textbf{M}_0,\textbf{M}_1\}$ un conjunto de matrices HPSD, siendo $\textbf{M}_0$ HPD, el producto matricial de Sobolev inducido, denotado por $\langle \cdot, \cdot \rangle_{\textbf{M}_{\textbf{S}}}$, se define como:
	\begin{equation}
		\langle p(t), q(t) \rangle_{\textbf{M}_{\textbf{S}}} = \langle p(t), q(t) \rangle_{\textbf{M}_0} + \langle p'(t), q'(t) \rangle_{\textbf{M}_1},\quad \forall p(t), q(t) \in \mathbb{P}[t]
	\end{equation}
\end{definition}

\begin{definition}\label{def:d_application}
	Sea $\mathcal{D}$ la aplicación $\mathcal{D}:\mathbb{P}[z] \mapsto \mathbb{P}[z]$ tal que: $D(p(z))=zp(z)$.
	Esta aplicación es lineal ya que
	$$\mathcal{D}(\alpha p(z) + \beta q(z))=z(\alpha p(z) + \beta q(z))=\alpha \mathcal{D}(p(z)) + \beta \mathcal{D}(q(z))$$
	\noindent Nos referiremos como $\mathbf{D}$ a la representación matricial del operador multiplicación por $z$ en la base de los polinomios ortonormales.
\end{definition}

\begin{definition}\label{def:bounded_eval_point}
	\cite{escribano2025zeros}
	Decimos que \(a\in\mathbb{C}\) es un \textit{punto de evaluación acotado} de una medida \(\mu\) (abreviado \emph{bpe} de \(\mu\)) si existe una constante \(C>0\) tal que
	\[
		|p(a)|^{2} \leq C \int |p(z)|^{2}\,d\mu(z), \qquad p(z)\in\mathbb{P}[z].
	\]
	Denotamos el conjunto de todos los puntos de evaluación acotados de \(\mu\) por
	\[
		bpe(\mu) =
		\left\{ a\in\mathbb{C} :\, a\ \text{es un bpe de } \mu \right\}.
	\]
\end{definition}

\begin{definition}\label{def:poly_convex_hull}
	\cite{escribano2025zeros}
	Para un conjunto compacto \(K \subset \mathbb{C}\), la \textit{envoltura convexa polinómica} de \(K\), denotada por \(P_C(K)\), se define como
	\[
		P_C(K) =
		\left\{ z \in \mathbb{C} :\, |p(z)| \leq \max_{\xi\in K}|p(\xi)|\ \text{ para todo } p\in\mathbb{P}[z] \right\}.
	\]
\end{definition}
Nótese que para el caso real, $P_C(K)$ coincide con los intervalos del soporte, y para el caso circunferencia, $P_C(K)$ coincide con las circunferencias del soporte.

\section{Estado del arte}

En la literatura se ha abordado el problema de los ceros de polinomios ortogonales de Sobolev mayoritariamente desde la perspectiva de la localización general de los ceros.

Primero vamos a probar que la acotacion del operador multiplicación $\mathcal{D}$ en $\mathbb{P}[z]$ es condicion suficiente para la acotación de los ceros de los polinomios ortogonales respecto a un producto inducido por una matriz HPD.

\begin{proposition}\label{prop:bounded_zeros}
	\cite{escribano2025boundedness} Sea \textbf{M} una matriz HPD y $\{\varphi_n(z)\}_{n=0}^{\infty}$ la secuencia de polinomios ortonormales asociados al producto interno inducido por \textbf{M}. Si $\mathcal{D}$ está acotado en $\mathbb{P}[z]$, el conjunto de ceros de los polonomios ortonormales asociados a \textbf{M} está acotado. Además,
	$$Z=\{z\in\mathbb{C}: \exists n \in \mathbb{N},\varphi_n(z)=0\} \subset \{z\in\mathbb{C}:|z|<\|\mathcal{D}\|\}.$$
\end{proposition}

\begin{proof}
	Asumamos que para $z_0 \in Z$ existe un $n \in \mathbb{N}$ tal que $\varphi_n(z_0)=0$. Luego, $\varphi_n(z)=(z-z_0)P_{n-1}(z)$, $0 \neq P_{n-1} \in \mathbb{P}_{n-1}[z]$. Por lo tanto, $\varphi_n(z_0)=zP_{n-1}(z)-z_0P_{n-1}(z)$ y porque $\varphi_n(z)$ y $P_{n-1}(z)$ son ortogonales en \textit{M} se sigue que
	$$\|zP_{n-1}(z)\|^2=\langle\varphi_n(z)+z_0P_{n-1}(z),\varphi_n(z)+z_0P_{n-1}(z)\rangle = \|\varphi_n(z)\|^2+|z_0|^2\|P_{n-1}(z)\|^2,$$
	por lo que,
	$$|z_0|^2\|P_{n-1}(z)\|^2=\|zP_{n-1}(z)\|^2- \|\varphi_n(z)\|^2<\|zP_{n-1}(z)\|^2=$$
	$$\|zP_{n-1}(z)\|^2=\|\mathcal{D}(P_{n-1})\|^2\leq\|\mathcal{D}\|^2\|P_{n-1}(z)\|^2.$$
	Como $\|P_{n-1}(z)\| \neq 0$, se concluye que
	$$|z_0|^2<\|\mathcal{D}\|,$$
	como se requería.
\end{proof}

Este resultado fundamental y usado extensamente en la literatura establece que los ceros de polinomios ortogonales de Sobolev están contenidos en un disco centrado en cero cuyo radio es la norma del operador $\mathcal{D}$. Veamos resultados que nos permiten acotar la norma del operador multiplicación $\mathcal{D}$.

Para acotar el operador multiplicación, G. López y H. Pijeira \cite{lagomasino1999zero} introdujeron el concepto de medidas secuencialmente dominadas. Diremos que el producto de Sobolev con dos medidas es secuencialmente dominado si
$$supp(\mu_2) \subset supp(\mu_{1}) , \qquad j = 1,2.$$ y
$$d\mu_2 = f_{1}d\mu_{1}, \qquad f_{1} \in L_{\infty}(\mu_{1}).$$
En nuestros casos continuos sería,
\begin{equation}
	[c,d] \subset [a,b] \quad y \quad\exists c > 0, w_1(t)\leq c \cdot w_0(t)
\end{equation}
y
\begin{equation}
	S_{a}(r_a)\subset S_{b}(r_b) ,\quad \text{la otra condición se tiene ya que},\quad w_1(t)\equiv w_2(t) \equiv 1
\end{equation}

Para trabajar con matrices de momentos introduciremos la definición para esta noción y un teorema análogo que nos dará la acotación.

\begin{definition}\label{def:seq_dominated}
	\cite{escribano2025boundedness} Sea $\{\textbf{M}_1,\textbf{M}_2\}$ un conjunto de matrices HSPD. Diremos que $\{\textbf{M}_1,\textbf{M}_2\}$ es secuencialmente dominada si existe una constante positiva $C$ tal que
	$$\textbf{M}_2 \leq C \cdot \textbf{M}_{1}.$$
\end{definition}

\begin{theorem}\label{thm:bounded_operator}
	\cite{escribano2025boundedness} Sea $\{\textbf{M}_1,\textbf{M}_2\}$ un conjunto de matrices HPD secuencialmente dominada. Asuma que para $j=1,2$ el operador multiplicación $\mathcal{D}_j$ está acotado en $(\mathbb{P}[z],\langle \cdot,\cdot \rangle)_{\textbf{M}_j}$. Considera el producto matricial de Sobolev siguiente:
	$$\langle p(t),q(t) \rangle_{\textbf{M}_S}=\langle p(t),q(t) \rangle_{\textbf{M}_1}+\langle p'(t),q'(t) \rangle_{\textbf{M}_2}.$$
	Entonces, el operador multiplicación $\mathcal{D}$ está acotado en $(\mathbb{P}[z],\langle \cdot,\cdot \rangle)_{\textbf{M}_S}.$
\end{theorem}

\begin{proposition}\label{prop:bounded_operator_continuous}
	Para el caso continuo real secuencialmente dominado con soportes $I_1=I_2=[a,b]$, definamos
	\(
	R = \max\{|a|,|b|\},
	\qquad
	\alpha = \max_{t\in[a,b]}\frac{w_1(t)}{w_0(t)}.
	\)
	Entonces el operador multiplicación $\mathcal{D}$
	es acotado en el producto Sobolev y satisface
	\(
	\|\mathcal{D}\|\;\le\;\sqrt{R^{2}+2\alpha}\,.
	\)
\end{proposition}

\begin{proof}
	Sea $p(t)\in\mathbb{P}[t]$.
	\noindent Utilizando que $(a+b)^2\leq 2(a^2+b^2)$ para $a,b\ge0$, se obtiene:
	$$
		\Vert \mathcal{D}(p(t))\Vert^{2}
		= \int_{a}^{b} t^2 p^2(t) w_0(t)\,dt
		\;+\;\int_{a}^{b} \Bigl[(tp(t))'\Bigr]^2 w_1(t)\,dt
	$$
	$$
		= \int_{a}^{b} t^2 p^2(t) w_0(t)\,dt
		+ \int_{a}^{b} \Bigl(p(t)+tp'(t)\Bigr)^2 w_1(t)\,dt
		\;\le\;
	$$
	$$
		\int_{a}^{b} t^2p^2(t)w_0(t)\; dt+2\Bigl(\int_{a}^{b} p(t)^2w_1(t)\; dt +\int_{a}^{b} t^2p'(t)^2w_1(t)\; dt\Bigr) \leq
	$$
	$$
		R^{2} \Bigl(\int_{a}^{b} p^2(t)w_0(t)\; dt +  2\int_{a}^{b} p'(t)^2w_1(t)\; dt\Bigr)+ 2\int_{a}^{b} p(t)^2w_1(t)\; dt \leq
	$$
	$$
		R^{2} \Vert p(t)\Vert^{2} + 2\alpha \int_{a}^{b} p(t)^2w_0(t)\; dt \leq (R^{2}+2\alpha)\Vert p(t) \Vert^{2}
	$$

\end{proof}

\noindent Por ejemplo, si $[a,b]=[-1,1]$ y $\alpha=1$, entonces
$\|\mathcal{D}\|\le\sqrt{1+2}=\sqrt{3}$.

Cuando $[c,d] \nsubseteq [a,b]$ entonces no podemos hablar de secuencialmente dominados. En este caso veremos el resultado de M. Alfaro y M. Luisa Rezola cuando estos soporte son disjuntos.

\begin{theorem}\label{thm:disjoint_support}
	\cite{garcia2001ceros} Si co($S_{\mu_0}$) $\cap$ co($S_{\mu_1}$) = $\emptyset$, entonces los ceros de $Q'_n$ son simples
	y están contenidos en el interior de co($S_{\mu_0}$) $\cup$ co($S_{\mu_1}$), y los ceros de $Q_n$ están
	contenidos en el disco centrado en el punto de co($S_{\mu_1}$) más alejado de co($S_{\mu_0}$) y
	radio el diámetro de co($S_{\mu_0}$) $\cup$ co($S_{\mu_1}$).
\end{theorem}

\begin{lemma}\label{lem:bounded_eval_point_in_hull}
	\cite{escribano2025zeros} Todo \emph{punto de evaluación acotado} de una medida $\mu$ está contenido en la envoltura polinómicamente convexa del soporte de $\mu$.

	\noindent Además, todo \emph{punto de evaluación acotado} de una medida $\mu$ es un punto de evaluación de la matriz de momentos asociada a $\mu$.
\end{lemma}

\begin{definition}\label{def:bounded_eval_point_matrix}
	\cite{escribano2025zeros} Sea $M$ una matriz hermítica positiva definida (HPD) y $a\in\mathbb{C}$.
	Decimos que $a$ es un \emph{punto de evaluación acotado} (bpe) de $M$ si existe una constante $C>0$ tal que
	\[ |p(a)|^{2}\;\le\;C\,\|p(z)\|_{M}^{2}, \quad p(z)\in\mathbb{P}[z]. \]
\end{definition}


\begin{proposition}\label{prop:equiv_bounded_eval_point}
	\cite{escribano2025zeros} Sea $M$ una matriz HPD y $a\in\mathbb{C}$. Entonces son equivalentes:
	\begin{enumerate}
		\item $a$ es un \emph{punto de evaluación acotado} de $M$;
		\item el conjunto $\{M,\,M(a)\}$ es secuencialmente dominado;
		\item $\gamma_{a}(M) > 0$.
	\end{enumerate}
\end{proposition}

\begin{proposition}\label{prop:bounded_derivative_eval}
	\cite{escribano2025zeros} Sea $M$ una matriz HPD tal que el operador de multiplicación $D$
	es acotado en $\bigl(\mathbb{P}[z],\|\cdot\|_{M}\bigr)$, y sean $a_{1},\dots,a_{n}\in\mathbb{C}$ evaluaciones puntuales acotadas de $M$.
	Entonces $D$ también es acotado en $\mathbb{P}[z]$ respecto a la norma
	\[
		\|p(z)\|_{MS}^{2}\;=\;\|p(z)\|_{M}^{2}\;+\;
		\sum_{k=1}^{n}\bigl|p'(a_{k})\bigr|^{2}.
	\]
\end{proposition}

\begin{corollary}\label{cor:bounded_zeros_circle}
	\cite{escribano2025zeros}
	Sea $m_{b;r_b}$ una medida soportada en una circunferencia $S_{b}(r_b)$ tal que los polinomios no son densos en $L^{2}(m_{b;r_b})$.
	Supongamos que $a_{1},\dots,a_{n}\in\operatorname{int}S_{b}(r_b)$ y consideremos la norma de Sobolev discreta
	\[
		\|p\|_{S}^{2}= \int_{S_{b}(r_b)} |p(z)|^{2}\,dm_{b;r_b}\;+\;\sum_{k=1}^{n} |p'(a_{k})|^{2},
		\qquad p\in\mathbb{P}[z].
	\]
	Entonces el operador multiplicación $D\colon p(z)\mapsto z\,p(z)$
	es acotado en $\bigl(\mathbb{P}[z],\|\cdot\|_{S}\bigr)$.
	En consecuencia, el conjunto de ceros de los polinomios ortogonales de Sobolev asociados está acotado.
\end{corollary}

\section{Motivación y Objetivos}

Como se ha puede observar, estudio de los ceros de los polinomios ortogonales de Sobolev se ha abordado desde el enfoque de acotar el operador multiplicación $\mathcal{D}$. En este trabajo nos proponemos estudiar el comportamiento asintótico de $\rho(\mathcal{D}_n)$, ya que sabemos que si este es acotado entonces todos los ceros lo serán. Para ello usaremos herramientas de cálculo númerico como Maple, que nos permitirá calcular y visualizar los polinomios ortogonales de Sobolev y sus ceros.

Siguiendo con el enfoque matricial propuesto por C. Escribano y R. Gonzalo, sabemos que los polinomios mónicos ortogonales de Sobolev se pueden obtener a partir de la matriz del operador multiplicación $\mathcal{D}$ y la matriz de momentos asociada al producto interno $\mathbf{M}$.
$$
	P_n(z)=\lvert zI_n-\mathcal{D}_n\rvert,\;
	\mathcal{D}=(c_{i,j})_{i,j=0}^\infty\;\text{con}\;
	c_{i,j}=\langle z\varphi_j,\varphi_i\rangle_{\mathbf{M}},\;
	\varphi_i,\varphi_j\;\text{ortonormales}
$$
Luego los ceros de los polinomios ortogonales de Sobolev se pueden obtener como los autovalores de la matriz $\mathcal{D}$. Por lo tanto, el estudio de los ceros de los polinomios ortogonales de Sobolev se reduce al estudio de los autovalores de la matriz $\mathcal{D}$, si esta fuese hermitiana sería diagonalizable con autovalores reales que coincidirían con los valores singulares, por la \textit{Definición \ref{def:singular_values}} sabemos que la norma del operador multiplicación $\mathcal{D}$ es igual al mayor valor singular $\sigma_{1}$, si este no fuese acotado entonces los ceros no estarían acotados. Sin embargo no podemos asegurar que la matriz $\mathcal{D}$ sea hermitiana, por lo que puede suceder que $\mathcal{D}$ no esta acotada y por el contrario los ceros de los polinomios ortogonales de Sobolev estén acotados, de pasar esto se intentará proponer otras condiciones de acotación utilizando la matriz $\mathcal{D}$.

